\section{Fixed masks and displacement moments}
\label{S:masked-moments}

We retain Hypothesis~\ref{H:homog}, data $(1,0)$, and all its original
dimensions and covariance amplitudes. Fix $\rho = (\epsilon,\eta) > 0$ and the
source $Y_\rho$ of Proposition~\ref{P:source-canonical-convergence}. For each
fixed $\rho$ use the continuous-time, smooth-space version in
Lemma~\ref{L:source-pointfield-regularity}. Source fields can use the auxiliary
extension preserving the given noise; their law and the given noise remain those
already constructed. Let $m \in C_c^\infty(\R^d)$ be real with $0 \le m \le 1$.
It multiplies the physical source after regularization.

Write $X_{\rho,n}^m(t,x)$ for the ordinary ordered displacement integral with
coefficient $mY_\rho$, and $X_{\rho,n}$ for its unmasked version; both
order-zero terms are one. We first construct these integrals. The positive Green
measures $G_d(r,\ud z)$ have mass $r$ and support in $B_r$, including the sphere
measure in $d = 3$. These are the source
kernels~\cite{chen:author2026:time-dependency}*{(2.1), (2.3)--(2.5), p.~7}; the
ordered convention and Gaussian pairing formula are
\cite{chen:author2026:time-dependency}*{(2.13)--(2.16), p.~10; (2.19), p.~12}.
We use all pair partitions, including self-pairings.

\begin{lemma}[Chains and mixed-order comparison]
  \label{L:masked-chain-comparison}
  The finite ordered integrals are Bochner $L^2$ integrals. For all $n,k \ge 0$,
  fixed $t,x,y$ and positive $\rho$,
  \begin{equation}\label{E:masked-order-comparison}
    0\le \E[X_{\rho,n}^m(t,x)X_{\rho,k}^m(t,y)]
    \le \E[X_{\rho,n}(t,x)X_{\rho,k}(t,y)].
  \end{equation}
  Both expectations vanish when $n+k$ is odd. With $q_n$ from
  Lemma~\ref{L:vertex-crude},
  \begin{equation}\label{E:masked-order-bound}
    \|X_{\rho,n}^m(t,x)\|_2\le q_n(t),\qquad
    \|X_{\rho,n}(t,x)\|_2\le q_n(t).
  \end{equation}
\end{lemma}

\begin{proof}
  We construct the ordered integrals using finite positive chain measures and
  Gaussian moment bounds. We then identify the Green-regularization limit and
  compare every Wick diagram before summing over pairings.

  \medskip
  \noindent\textit{Step 1.~}We construct the finite chain integrals.
  Let $\nu_n^{t,x}$ be the positive time-space chain measure on $0 < s_1 <
  \cdots < s_n < t$, with successive spatial factors $G_d(s_{j+1}-s_j,\ud z_j)$
  and $s_{n+1} = t$, recovering positions backwards from $x$. It uses measures
  rather than sphere densities. The order-zero measure has mass one. Integrating
  the latest vertex gives the recursion $f_n(t) = \int_0^t(t-r)f_{n-1}(r)\ud r$;
  hence,
  \begin{equation}\label{E:masked-chain-mass}
    \nu_n^{t,x}(1)=\frac{t^{2n}}{(2n)!}.
  \end{equation}
  In the covariance and mask factors below, $z_i$ denotes the vertex position
  recovered backwards from the terminal point, not the Green increment used to
  construct the chain measure.
  Starting from mass one at order zero, the recursion integrates the mass of
  the last Green factor, $t-r$, against the previous mass. Substituting
  $f_{n-1}(r) = r^{2n-2}/(2n-2)!$ gives the displayed mass by elementary
  integration. In particular, the measure is finite at each fixed order.
  Measurability follows by representing $G_d(r)$ as $r$ times the dilation by
  $r$ of its unit-time probability measure.

  The covariance is
  \[
    C_\rho(s,z;r,w)=K_\eta(s-r)\Gamma_{2\epsilon}(z-w),
    \qquad 0\le C_\rho\le c_\rho
    :=K_\eta(0)\Gamma_{2\epsilon}(0)<\infty.
  \]
  Positivity uses the nonnegative Radon measure $\Gamma$ and positive heat and
  Laplace kernels. The upper bound also follows by Gaussian Cauchy--Schwarz.
  H\"older's inequality and scalar Gaussian moments bound the absolute
  expectation of any fixed number of source factors, uniformly over their
  vertices. Applied to $2n$ factors this bounds the $L^2$ norm of each $n$-fold
  product. The field is measurable and continuous, so integration against the
  finite chain measure defines a Bochner $L^2$ integral and permits stochastic
  Fubini between two chains.

  The bound on the $L^2$ norm of each product is uniform over the integration
  variables at fixed $\rho,n$. Its integral against the finite chain measure
  is therefore finite. This is the Bochner integrability condition, and
  applying the same argument to two products gives the absolute integrability
  needed to exchange expectation and the two chain integrals.

  \medskip
  \noindent\textit{Step 2.~}We identify the integrals obtained by Green
  smoothing.
  These are the canonical regularized terminal integrals at common noise widths.
  The unmasked $X_{\rho,n}(t,x)$ is $Y_n^{\boldsymbol a}(t,x)$ from
  Lemma~\ref{L:vertex-crude} with every vertexwise noise width in $\boldsymbol
  a$ set equal to $\rho$, so its bound is $q_n(t)$.
  Here is the Green-width identification needed in dimension three. At fixed
  positive $\rho$, replace every Green measure by its convolution with a
  positive heat approximate identity. At fixed times the measures converge
  weakly, preserving masses. Every Wick covariance product is bounded by a power
  of $c_\rho$ and is spatially continuous by
  Lemma~\ref{L:source-pointfield-regularity}. Thus, its spatial integrals
  converge. The chain mass integrands from~\eqref{E:masked-chain-mass} give an
  integrable time majorant. Square and cross moments between smoothed and
  unsmoothed integrals converge to the same diagram sum; therefore, the random
  integrals converge in $L^2$. This is precisely the Green-before-noise limit in
  Lemma~\ref{L:joint-fixed-order}. It also applies in dimensions one and two if
  Green smoothing is used. No samplewise source bound on heat-kernel tails is
  needed. Consequently, Lemma~\ref{L:vertex-crude} applies to the unmasked
  integrals.


  \medskip
  \noindent\textit{Step 3.~}We compare each masked Wick diagram with its
  unmasked diagram.
  If $n+k$ is odd, Gaussian centering gives zero. If $n+k = 2l$, each pair
  partition $D$ contributes the nonnegative finite integral
  \[
    \int \prod_{\{i,j\}\in D}
    C_\rho(s_i,z_i;s_j,z_j)
    \ud(\nu_n^{t,x}\otimes\nu_k^{t,y}).
  \]
  The masked integrand has the additional factor $\prod_{i=1}^{2l}m(z_i) \in
  [0,1]$. Every vertex appears exactly once, irrespective of its chain or
  partner. Summing gives \eqref{E:masked-order-comparison}, including unequal
  and empty chains and every self-pairing. At $n = k = 0$ both sides are one.
  Taking $k = n$, $y = x$ and then applying the unmasked bound gives
  \eqref{E:masked-order-bound}. This is a statement about expected displacement
  diagrams; no samplewise solution order or Gaussian covariance order is used.
  This proves Lemma~\ref{L:masked-chain-comparison}.
\end{proof}

\begin{proposition}[Ordinary solutions and tested moments]
  \label{P:masked-source-moments}
  The sums $u_\rho^m = \sum_{n\ge0}X_{\rho,n}^m$ and $u_\rho =
  \sum_{n\ge0}X_{\rho,n}$ converge absolutely in $L^2$ and have locally
  continuous versions that are the unique locally continuous ordinary mild
  solutions with coefficients $mY_\rho$ and $Y_\rho$, respectively, and data
  $(1,0)$. For every nonnegative $\phi \in C_c^\infty(\R^d)$,
  \begin{align}
    0&\le\E[u_\rho^m(t,x)u_\rho^m(t,y)]
    \le\E[u_\rho(t,x)u_\rho(t,y)],\label{E:masked-two-point}\\
    \E\left|\int\phi(x)u_\rho^m(t,x)\ud x\right|^2
    &\le \|\phi\|_1^2\,\E|u_\rho(t,0)|^2.
    \label{E:masked-average-source}
  \end{align}
  The spatial integral is an Bochner $L^2$ integral.
\end{proposition}

\begin{proof}
  We first identify the mean-square series with the continuous Picard solution
  on each compact observation region. We then pass the mixed-order comparison
  through the series and integrate it against the nonnegative spatial test.

  \medskip
  \noindent\textit{Step 1.~}We construct and identify the ordinary mild
  solutions.
  Put $Q_T = \sum_{n\ge0}q_n(T) < \infty$, using~\eqref{E:vertex-summability}.
  Lemma~\ref{L:masked-chain-comparison} proves absolute $L^2$ convergence
  uniformly in space and widths. To identify the sums pathwise, fix $\rho$, a
  compact observation set and a finite horizon $T$. On the source regularity
  event the coefficient is bounded by a finite sample-dependent $M$ on the
  compact backward-cone enlargement. Finite propagation
  and~\eqref{E:masked-chain-mass} bound each Picard term there by
  $M^nT^{2n}/(2n)!$. This series converges uniformly. Each term is continuous by
  the unit-time Green representation and dominated integration. Summing the
  Volterra recursion gives the ordinary mild equation with initial displacement
  one and velocity zero. Iterating the difference equation between two
  continuous solutions bounds its value on the observation set by $B
  M^nT^{2n}/(2n)!$, where $B$ bounds their difference on the larger compact
  cone. Letting $n \to \infty$ proves uniqueness.


  \medskip
  \noindent\textit{Step 2.~}We identify the continuous version with the
  Bochner series.
  Ordinary Fubini and the absolute product integrability in the preceding lemma
  identify each Bochner term with the pathwise chain integral at a fixed
  terminal point, simultaneously over countably many orders. An almost surely
  convergent subsequence of the $L^2$ partial sums then identifies its limit
  with the continuous Picard version at that point. These versions define the
  solution for all terminal points at this fixed $\rho$; no common event over
  uncountably many widths is needed.

  \medskip
  \noindent\textit{Step 3.~}We sum the expected two-point comparison.
  For either masked or unmasked terms, Cauchy--Schwarz gives
  \[
    \sum_{n,k\ge0}\E|X_{\rho,n}^{(m)}(t,x)
    X_{\rho,k}^{(m)}(t,y)|\le Q_t^2.
  \]
  For the unmasked stationary field, Cauchy--Schwarz bounds its two-point
  expectation by the geometric mean of two equal point second moments.
  Products of partial sums converge in $L^1$. Thus, the preceding lemma passes
  termwise to~\eqref{E:masked-two-point}. The unmasked source law is spatially
  stationary, and its finite Picard integrals commute with translations. Their
  $L^2$ limits inherit stationarity, so $\E[u_\rho(t,x)u_\rho(t,y)] \le
  \E|u_\rho(t,0)|^2$. Stationarity is not required of the masked solution.


  \medskip
  \noindent\textit{Step 4.~}We form the tested spatial average.
  Translating every recovered vertex position by $x$ identifies $\nu_n^{t,x}$
  with the translate of $\nu_n^{t,0}$; on this fixed increment measure,
  varying $x$ changes only the arguments of the source factors and, in the
  masked case, of $m$.
  Each finite Bochner term is $L^2$-continuous in $x$ by source increment
  continuity on compact cones and dominated Gaussian product moments. Uniform
  $q_n$ tails pass this property to the full sum. Its separable $L^2$ image and
  bound $Q_t$ give the Bochner average. Absolute Fubini is justified by
  $\|\phi\|_1^2Q_t^2$. Integrating the two-point inequality against
  $\phi(x)\phi(y)$ proves~\eqref{E:masked-average-source}. This proves
  Proposition~\ref{P:masked-source-moments}.
\end{proof}

\begin{proposition}[Canonical point-moment limit and conditional transfer]
  \label{P:masked-canonical-transfer}
  At every fixed $t$, at arbitrary relative rates of the two widths,
  \begin{equation}\label{E:masked-unmasked-moment-limit}
    \lim_{\epsilon,\eta\downarrow0}\E|u_\rho(t,0)|^2
    =\E|u(t,0)|^2,
  \end{equation}
  where $u$ is the original canonical Stratonovich solution. If the masked
  averages in~\eqref{E:masked-average-source} converge in $L^2$ to
  $Z_m(t,\phi)$, either at arbitrary rates or along a specified sequence with
  both widths tending to zero, then
  \begin{equation}\label{E:masked-conditional-transfer}
    \E|Z_m(t,\phi)|^2\le\|\phi\|_1^2\,\E|u(t,0)|^2.
  \end{equation}
\end{proposition}

\begin{proof}
  We first pass to the canonical limit at each pair of chain orders. A
  summable double majorant then gives the full point moment. The conditional
  masked transfer follows by passing to the limit in the tested source
  inequality.

  \medskip
  \noindent\textit{Step 1.~}We identify the limit for fixed orders.
  The relabelling $r$ below reverses the vertex order separately within each
  chain.
  At every fixed pair of orders the common-width source law is the vertexwise
  law in Lemma~\ref{L:joint-fixed-order}. At terminal $(t,0)$ reverse each chain
  by $s'_j = t-s_{n+1-j}$, $z'_j = -z_{n+1-j}$, and likewise for the other
  chain. Even covariances preserve the reversed differences. Relabeling vertices
  sends each pairing $\{i,j\}$ to $\{r(i),r(j)\}$, giving the exact
  initial-chain convention in Proposition~\ref{P:mixed}, also for an empty
  chain. The preceding Green-width identification removes Green smoothing first.
  Lemma~\ref{L:joint-fixed-order} then gives the canonical limit of each diagram
  at arbitrary relative noise-width rates. Odd mixed moments stay zero. Hence,
  at fixed $n,k$,
  \[
    \E[X_{\rho,n}(t,0)X_{\rho,k}(t,0)]
    \longrightarrow \E[X_n(t,0)X_k(t,0)].
  \]

  \medskip
  \noindent\textit{Step 2.~}We pass from finitely many orders to the full
  moment.
  All absolute mixed moments, before and after the limit, are bounded by
  $q_n(t)q_k(t)$. Given $\delta > 0$, choose a finite square of orders with the
  double majorant outside it less than $\delta$. On the square the finitely many
  limits hold jointly; off it both tails are bounded by $\delta$.
  Proposition~\ref{P:mixed} identifies the canonical sum, proving
  \eqref{E:masked-unmasked-moment-limit}. This argument uses the single-width
  diagram laws, without assigning a new pathwise realization of the entire
  canonical solution family.

  \medskip
  \noindent\textit{Step 3.~}We transfer the bound under the stated
  masked-limit premise.
  Finally, $L^2$ convergence of the masked averages implies convergence of their
  second moments by
  \[
    |\E|B|^2-\E|A|^2|
    \le\|B-A\|_2(\|B\|_2+\|A\|_2).
  \]
  The inequality follows by factoring the difference of squared norms and
  applying Cauchy--Schwarz. Convergence in $L^2$ also bounds both norms on the
  right, so this difference tends to zero. Thus, both sides of the source
  inequality have the limits needed for the comparison.
  Take the same source limit in~\eqref{E:masked-average-source} and
  use~\eqref{E:masked-unmasked-moment-limit}. Existence of the masked limit
  remains an explicit premise here; the next section constructs it on a phase
  containing the constant data. This proves
  Proposition~\ref{P:masked-canonical-transfer}.
\end{proof}
